\documentclass[10pt]{article}
\usepackage[margin=0.8in]{geometry}
\usepackage{amsmath,amssymb,booktabs}
\begin{document}
\section*{Conditional fraction-field pole identity}
The pinned Lean file \texttt{Common22FractionPole.lean} kernel-checks the
fraction-field valuation step conditional on the infinity chart.  Let $v$
be a valuation on a field, and assume $w\ne0$, $w^2=zu$, $v(u)=1$, and
$z^3g=G$. The two checked implications are:
\begin{center}
\begin{tabular}{lll}\toprule
Theorem & Additional hypothesis & Conclusion\\\midrule
\texttt{unit\_chart\_pole} & $v(G)=1$ & $v(g)=v(w)^{-6}$\\
\texttt{linear\_chart\_pole} & $v(G)=v(w)$ & $v(g)=v(w)^{-5}$\\\bottomrule
\end{tabular}
\end{center}
The proof uses multiplicativity of $v$, the chart relation, and cancellation
of the nonzero value $v(w)$. The first branch corresponds to a unit chart
residue; the second to the degree-drop branch. These exponents become pole
orders six and five only after $w$ is identified as a normalized local
uniformizer of the actual smooth projective curve.

\paragraph{Boundary.} The projective infinity stalk, its embedding in the
function field, the principal divisor, Mumford reduction, the $\operatorname{Pic}^0$
class equation, and compiled Rust semantics are not supplied by this file.
The result changes no target solve or cost ratio. The adjacent CSV records
both theorem statuses and a null rho ratio; the SVG graph displays only
the conditional exponents.

\paragraph{Replay.} Run \texttt{lake env lean Common22FractionPole.lean} in
the pinned formal directory. The successful axiom output is retained in
\texttt{fraction-pole-build.log}.
\end{document}
